\section{引言}
\label{sec:intro}

零样本异常检测的目标是：在目标产品或域上没有任何训练样本（正常与异常皆无）的前提下，判定图像或像素是否异常。CLIP 因其开放词表的视觉-文本对齐能力成为该任务的主流基座~\cite{radford2021clip}：WinCLIP~\cite{jeong2023winclip} 与 AnomalyCLIP~\cite{zhou2024anomalyclip} 等方法用文本原型描述“正常/异常”两种状态，再与图像特征比对得到异常分数。

然而 CLIP 存在一个结构性、而非调参可解决的缺陷：它被训练用于捕捉物体级的\emph{全局}语义，其 patch token 经注意力聚合后偏\emph{低通}。而工业缺陷往往不改变物体类别，只体现为局部纹理或高频的细微变化。既有的 CLIP-ZSAD 工作大多在\emph{语义空间}里修补这一缺陷——WinCLIP 用多尺度窗口聚合，AnomalyCLIP 学习物体无关的文本原型，VCP-CLIP~\cite{qu2024vcpclip} 与 AA-CLIP~\cite{ma2025aaclip} 在特征空间做适配；近期的频率类工作~\cite{gong2025feclip} 则把频率作为\emph{额外特征、经可训练适配器融入} CLIP。

这些工作都默认“CLIP 缺少异常线索、需要补充”。本文提出一个不同的问题：异常线索是否其实已经存在于 CLIP 中，只是无法直接使用？我们的分析给出肯定回答，并进一步指出它为什么不能直接用，以及用什么最小的、训练无关的步骤能让它变得可用。围绕三个问题展开：（Q1）异常相关的高频信号是否已经存在于 CLIP patch 特征中？（Q2）若存在，为什么不能直接使用？（Q3）什么样的最小、训练无关的步骤能让它变得可用？

\figref{fig:motivation} 给出核心直觉。对于光滑物体上的缺陷（左），高频响应能干净地把缺陷凸显出来；但对于正常的粗糙材质（右），\emph{整片}区域的高频都很高。二者的高频幅度可能相同，含义却相反——因此任何固定阈值都必然失败。缺的正是“这张图自己的正常高频水平”这一参照。

\begin{figure*}[t]
\centering
\begin{tikzpicture}[
  font=\small,
  box/.style={rounded corners=2pt,draw=black!30,thick},
  lbl/.style={font=\footnotesize\bfseries,text=cInk},
]
% ===== 左：光滑物体+缺陷 =====
\node[lbl] at (1.6,3.3) {(a) 光滑物体 + 缺陷};
% input
\node[box,fill=cBase!35,minimum width=2.4cm,minimum height=1.5cm] (ia) at (1.6,2.2) {};
\fill[cBad] (1.35,2.35) rectangle (2.05,2.5);
\node[font=\scriptsize,text=black!55] at (1.6,1.25) {输入};
% HF response
\node[box,fill=cFreq!22,minimum width=2.4cm,minimum height=1.5cm] (ha) at (1.6,-0.6) {};
\fill[cFreq] (1.35,-0.45) rectangle (2.05,-0.3);
\node[font=\scriptsize,text=cFreq!60!black] at (1.6,-1.55) {缺陷凸显（好）};

% ===== 右：粗糙正常纹理 =====
\node[lbl] at (6.0,3.3) {(b) 粗糙材质，正常};
\node[box,fill=cBase!35,minimum width=2.4cm,minimum height=1.5cm] (ib) at (6.0,2.2) {};
\foreach \y in {1.65,1.85,2.05,2.25,2.45,2.65}{\draw[black!35,line width=0.5pt] (4.85,\y)--(7.15,\y);}
\node[font=\scriptsize,text=black!55] at (6.0,1.25) {无缺陷};
\node[box,fill=cFreq!22,minimum width=2.4cm,minimum height=1.5cm] (hb) at (6.0,-0.6) {};
\foreach \y in {-1.15,-0.95,-0.75,-0.55,-0.35,-0.15}{\draw[cFreq!80!black,line width=0.6pt,opacity=0.75] (4.85,\y)--(7.15,\y);}
\node[font=\scriptsize,text=cGlobal!45!black] at (6.0,-1.55) {整片变亮（坏）};

% 行标题
\node[lbl,rotate=90,text=cFreq!60!black] at (-0.35,-0.6) {CLIP 高频响应};
\node[lbl,rotate=90,text=black!55] at (-0.35,2.2) {输入 / 真值};

% 箭头
\draw[-{Latex[length=2mm]},black!55] (ia) -- (ha);
\draw[-{Latex[length=2mm]},black!55] (ib) -- (hb);

% ===== 中间洞察框 =====
\node[box,fill=cGlobal!18,draw=cGlobal!70,minimum width=3.2cm,align=center,font=\footnotesize] (ins) at (9.9,1.7)
{\textbf{相同的高频幅度}\\[1pt] 在 (a) 是异常\\ 在 (b) 是正常\\[2pt]\textbf{$\Rightarrow$ 固定阈值失效}};

% ===== 解法框 =====
\node[box,fill=cOurs!18,draw=cOurs!70,minimum width=6.6cm,align=center,font=\footnotesize] (sol) at (11.6,-0.9)
{\textbf{本文思路：} 估计每张图\emph{自己}的正常高频水平，只标记超出该逐图参照的部分};

\draw[-{Latex[length=2mm]},cOurs!70,thick] (ins) -- (sol);
\end{tikzpicture}
\caption{CLIP patch 特征的高频响应对缺陷（a）敏感，但对正常粗糙材质（b）同样敏感。单一固定阈值无法区分“异常”与“材质”。我们的方法估计每张图自己的正常高频水平，只标记超出该逐图参照的部分。}
\label{fig:motivation}
\end{figure*}


\paragraph{贡献。}
\begin{itemize}[leftmargin=1.4em,itemsep=0.15em,topsep=0.2em]
\item \textbf{一个诊断。} CLIP patch 特征已携带异常相关的高频信号，但因正常材质同样高频而\emph{歧义}；我们以受控消融量化之：直接使用会导致异常图崩盘。
\item \textbf{一个训练无关的方法。} 读出高频，用这个 CLIP 之外的信号选择证据，并\emph{逐图估计正常参照}；异常即相对该参照的偏离。
\item \textbf{机制证据。} 参照必须（a）逐图而非全局、（b）由独立于 CLIP 语义的信号驱动；我们召回约 $18\text{--}22\%$ 被 CLIP 独自漏检的异常，且增益集中在频率盲的纹理类别。
\end{itemize}

我们在五个数据集上取得对原始 AnomalyCLIP 的一致提升，且核心消融呈现清晰的单调序：DirectHF $<$ Baseline $<$ GlobalRef $<$ SelfRef $<$ Ours。
